As eqs. (4.53) e (4.54) representam, respectivamente, a \textbf{inclinação} $m$ e o \textbf{intercepto} $b$ da reta ajustada pelo método dos mínimos quadrados. As suas dimensões são distintas:

\begin{itemize}
\item \textbf{Eq. (4.53) --- inclinação} $m$:
\[
m = \frac{n\sum x_i y_i - \sum x_i\sum y_i}{n\sum x_i^{2}-(\sum x_i)^{2}} \implies [m] = \frac{[y]}{[x]}.
\]
\item \textbf{Eq. (4.54) --- intercepto} $b$:
\[
b = \frac{\sum y_i - m\sum x_i}{n} \implies [b] = [y] - \frac{[y]}{[x]}\cdot[x] = [y].
\]
\end{itemize}

A diferença dimensional essencial é: $m$ tem as dimensões de \textbf{razão} entre as variáveis ($y/x$), enquanto $b$ tem as dimensões da \textbf{própria variável dependente} ($y$). As duas equações não podem ser somadas ou comparadas diretamente; cada uma contribui com um papel físico diferente na equação da reta $y = mx + b$. 